\section*{Problemas de regresión lineal múltiple}

\subsection{Enunciados de los problemas}

\begin{problema}
	Se estudia el rendimiento de combustible ($Y$, en km/l) de automóviles en función de dos variables predictoras: el peso del vehículo ($X_1$, en toneladas) y la cilindrada del motor ($X_2$, en litros). Se ajustó un modelo de regresión múltiple con $n = 30$ observaciones y se obtuvieron los siguientes resultados:
	
	\begin{align}
		\hat{Y} = 25 - 3X_1 - 1.5X_2, \quad R^2 = 0.78, \quad R^2_{adj} = 0.76
	\end{align}
	
	Los errores estándar de los coeficientes son: $SE(\hat{\beta}_0) = 2.5$, $SE(\hat{\beta}_1) = 0.8$, $SE(\hat{\beta}_2) = 0.6$.
	
	\begin{enumerate}
		\item Interprete los coeficientes $\hat{\beta}_1$ y $\hat{\beta}_2$.
		\item Prediga el rendimiento de un automóvil que pesa 1.5 toneladas y tiene motor de 2.0 litros.
		\item Realice pruebas de hipótesis individuales para cada coeficiente con $\alpha = 0.05$.
		\item ¿Qué variable predictora tiene mayor impacto en el rendimiento?
	\end{enumerate}
\end{problema}

\begin{problema}
	Una empresa quiere predecir las ventas trimestrales ($Y$, en miles de dólares) usando tres predictores: gasto en publicidad en TV ($X_1$), gasto en publicidad en radio ($X_2$) y número de vendedores ($X_3$). Con $n = 40$ trimestres, se obtuvieron dos modelos:
	
	\textbf{Modelo 1:} $\hat{Y} = 50 + 2.5X_1 + 1.8X_2$, con $R^2 = 0.75$
	
	\textbf{Modelo 2:} $\hat{Y} = 45 + 2.3X_1 + 1.6X_2 + 0.8X_3$, con $R^2 = 0.78$
	
	\begin{enumerate}
		\item Calcule el $R^2$ ajustado para ambos modelos.
		\item Realice una prueba F parcial para determinar si agregar $X_3$ mejora significativamente el modelo.
		\item ¿Qué modelo recomendaría y por qué?
	\end{enumerate}
\end{problema}

\begin{problema}
	Se ajustó un modelo de regresión múltiple para predecir el precio de viviendas ($Y$, en miles de dólares) usando: área ($X_1$, en m²), número de habitaciones ($X_2$), antigüedad ($X_3$, en años) y distancia al centro ($X_4$, en km). Con $n = 50$ observaciones:
	
	\begin{align}
		\hat{Y} = 50 + 0.8X_1 + 15X_2 - 2X_3 - 5X_4
	\end{align}
	
	Los valores VIF son: $VIF(X_1) = 4.2$, $VIF(X_2) = 3.8$, $VIF(X_3) = 1.5$, $VIF(X_4) = 1.2$.
	
	\begin{enumerate}
		\item Interprete los coeficientes de $X_1$ y $X_3$.
		\item ¿Hay evidencia de multicolinealidad? ¿Qué variables están involucradas?
		\item Prediga el precio de una casa con 120 m², 3 habitaciones, 10 años de antigüedad y a 8 km del centro.
		\item ¿Qué acción tomaría respecto a la multicolinealidad?
	\end{enumerate}
\end{problema}

\begin{problema}
	Un investigador estudia el salario ($Y$, en miles de dólares) en función de años de experiencia ($X_1$) y nivel educativo ($X_2$, codificado como: 1=bachillerato, 2=licenciatura, 3=maestría, 4=doctorado). Con $n = 100$ empleados:
	
	\begin{align}
		\hat{Y} = 20 + 2X_1 + 8X_2, \quad R^2 = 0.65
	\end{align}
	
	Errores estándar: $SE(\hat{\beta}_1) = 0.3$, $SE(\hat{\beta}_2) = 1.5$
	
	\begin{enumerate}
		\item Interprete el coeficiente $\hat{\beta}_2$.
		\item ¿Es apropiado tratar $X_2$ como variable cuantitativa? ¿Qué alternativa propondría?
		\item Prediga el salario de alguien con 10 años de experiencia y maestría.
		\item Construya un intervalo de confianza del 95\% para $\beta_1$.
	\end{enumerate}
\end{problema}

\begin{problema}
	Se quiere predecir el tiempo de entrega ($Y$, en días) usando: distancia ($X_1$, en km), peso del paquete ($X_2$, en kg) y tipo de servicio ($X_3$, donde 1=express, 0=estándar). Con $n = 60$ entregas:
	
	\textbf{Modelo completo:} $\hat{Y} = 2 + 0.01X_1 + 0.5X_2 - 3X_3$, con $R^2 = 0.82$, $SSD = 45$
	
	\textbf{Modelo reducido (sin $X_3$):} $\hat{Y} = 3.5 + 0.012X_1 + 0.55X_2$, con $R^2 = 0.68$, $SSD = 78$
	
	\begin{enumerate}
		\item Interprete el coeficiente de $X_3$ en el modelo completo.
		\item Realice una prueba F parcial para determinar si $X_3$ es significativo.
		\item Calcule el $R^2$ ajustado para ambos modelos.
		\item ¿Vale la pena incluir el tipo de servicio en el modelo?
	\end{enumerate}
\end{problema}

\begin{problema}
	Se ajusta un modelo para predecir el rendimiento académico ($Y$, promedio de 0 a 10) usando horas de estudio semanales ($X_1$), asistencia a clase ($X_2$, porcentaje de 0 a 100) y uso de recursos en línea ($X_3$, horas semanales). Con $n = 80$ estudiantes:
	
	\begin{align}
		\hat{Y} = 4 + 0.15X_1 + 0.03X_2 + 0.1X_3
	\end{align}
	
	Estadísticos: $R^2 = 0.58$, $F = 35.2$ (con $p$-valor $< 0.001$)
	
	Valores $p$ individuales: $p_1 < 0.001$, $p_2 = 0.023$, $p_3 = 0.145$
	
	\begin{enumerate}
		\item ¿Es significativo el modelo global?
		\item ¿Qué variables son significativas individualmente con $\alpha = 0.05$?
		\item ¿Qué acción tomaría respecto a $X_3$?
		\item Si elimina $X_3$, ¿qué esperaría que suceda con $R^2$ y $R^2_{adj}$?
	\end{enumerate}
\end{problema}


\subsection{Soluciones detalladas}

\begin{solucion}
	\begin{enumerate}
		\item \textbf{Interpretación:}
		\begin{itemize}
			\item $\hat{\beta}_1 = -3$: Manteniendo constante la cilindrada, por cada tonelada adicional de peso, el rendimiento disminuye en promedio 3 km/l.
			\item $\hat{\beta}_2 = -1.5$: Manteniendo constante el peso, por cada litro adicional de cilindrada, el rendimiento disminuye en promedio 1.5 km/l.
		\end{itemize}
		
		\item Predicción para $X_1 = 1.5$, $X_2 = 2.0$:
		\begin{align}
			\hat{Y} = 25 - 3(1.5) - 1.5(2.0) = 25 - 4.5 - 3.0 = 17.5 \text{ km/l}
		\end{align}
		
		\item Pruebas de hipótesis ($H_0: \beta_i = 0$ vs $H_1: \beta_i \neq 0$):
		
		Para $\beta_1$:
		\begin{align}
			t_1 = \frac{-3}{0.8} = -3.75, \quad |t_1| = 3.75
		\end{align}
		
		Para $\beta_2$:
		\begin{align}
			t_2 = \frac{-1.5}{0.6} = -2.5, \quad |t_2| = 2.5
		\end{align}
		
		Valor crítico: $t_{0.025, 27} \approx 2.052$
		
		Ambos estadísticos exceden el valor crítico, por lo que rechazamos $H_0$ para ambos coeficientes. Tanto el peso como la cilindrada son predictores significativos al nivel $\alpha = 0.05$.
		
		\item Para comparar el impacto relativo, usamos coeficientes estandarizados. El peso tiene un impacto mayor (coeficiente absoluto de 3 vs 1.5), aunque esto también depende de las escalas de las variables.
	\end{enumerate}
\end{solucion}

\begin{solucion}
	\begin{enumerate}
		\item $R^2$ ajustado:
		
		Modelo 1 ($p = 2$):
		\begin{align}
			R^2_{adj,1} = 1 - \frac{40-1}{40-2-1}(1 - 0.75) = 1 - \frac{39}{37}(0.25) = 1 - 0.264 = 0.736
		\end{align}
		
		Modelo 2 ($p = 3$):
		\begin{align}
			R^2_{adj,2} = 1 - \frac{40-1}{40-3-1}(1 - 0.78) = 1 - \frac{39}{36}(0.22) = 1 - 0.238 = 0.762
		\end{align}
		
		\item Prueba F parcial para $X_3$:
		\begin{align}
			F &= \frac{(R^2_2 - R^2_1)/(p_2 - p_1)}{(1 - R^2_2)/(n - p_2 - 1)} = \frac{(0.78 - 0.75)/1}{(1 - 0.78)/(40 - 3 - 1)} \\
			&= \frac{0.03}{0.22/36} = \frac{0.03}{0.00611} \approx 4.91
		\end{align}
		
		Valor crítico: $F_{0.05, 1, 36} \approx 4.11$
		
		Como $F = 4.91 > 4.11$, rechazamos $H_0$. Agregar $X_3$ mejora significativamente el modelo al nivel $\alpha = 0.05$.
		
		\item \textbf{Recomendación:} El Modelo 2 es preferible porque:
		\begin{itemize}
			\item Tiene mayor $R^2$ ajustado (0.762 vs 0.736)
			\item La variable $X_3$ es estadísticamente significativa
			\item La mejora en $R^2$ es significativa según la prueba F parcial
		\end{itemize}
	\end{enumerate}
\end{solucion}

\begin{solucion}
	\begin{enumerate}
		\item \textbf{Interpretación:}
		\begin{itemize}
			\item $\hat{\beta}_1 = 0.8$: Manteniendo constantes las demás variables, por cada m² adicional de área, el precio aumenta en promedio 800 dólares.
			\item $\hat{\beta}_3 = -2$: Manteniendo constantes las demás variables, por cada año adicional de antigüedad, el precio disminuye en promedio 2,000 dólares.
		\end{itemize}
		
		\item \textbf{Análisis de multicolinealidad:}
		\begin{itemize}
			\item $VIF(X_1) = 4.2$: Moderada (cercano al umbral de 5)
			\item $VIF(X_2) = 3.8$: Moderada
			\item $VIF(X_3) = 1.5$: Baja
			\item $VIF(X_4) = 1.2$: Baja
		\end{itemize}
		
		Hay evidencia de multicolinealidad moderada entre $X_1$ (área) y $X_2$ (habitaciones), lo cual es lógico: casas más grandes tienden a tener más habitaciones.
		
		\item Predicción para $X_1 = 120$, $X_2 = 3$, $X_3 = 10$, $X_4 = 8$:
		\begin{align}
			\hat{Y} &= 50 + 0.8(120) + 15(3) - 2(10) - 5(8) \\
			&= 50 + 96 + 45 - 20 - 40 = 131
		\end{align}
		El precio predicho es 131,000 dólares.
		
		\item \textbf{Acciones recomendadas:}
		\begin{itemize}
			\item Monitorear los errores estándar de $\hat{\beta}_1$ y $\hat{\beta}_2$
			\item Considerar eliminar una de las dos variables si los errores estándar son muy grandes
			\item Alternativamente, crear una variable compuesta (ej: área por habitación)
			\item Si el modelo tiene buen poder predictivo, la multicolinealidad moderada puede ser aceptable
		\end{itemize}
	\end{enumerate}
\end{solucion}

\begin{solucion}
	\begin{enumerate}
		\item \textbf{Interpretación de $\hat{\beta}_2$:} Manteniendo constantes los años de experiencia, por cada nivel educativo superior, el salario aumenta en promedio 8,000 dólares. Sin embargo, esta interpretación asume que la diferencia entre niveles es constante, lo cual puede no ser realista.
		
		\item \textbf{Apropiación del tratamiento:} No es completamente apropiado tratar $X_2$ como cuantitativa porque:
		\begin{itemize}
			\item Los niveles no están equiespaciados (la diferencia entre bachillerato y licenciatura no es necesariamente igual a la diferencia entre maestría y doctorado)
			\item Se asume una relación lineal que puede no existir
		\end{itemize}
		
		\textbf{Alternativa:} Usar variables dummy (categóricas) para cada nivel educativo, tomando uno como referencia (ej: bachillerato).
		
		\item Predicción para $X_1 = 10$, $X_2 = 3$ (maestría):
		\begin{align}
			\hat{Y} = 20 + 2(10) + 8(3) = 20 + 20 + 24 = 64
		\end{align}
		El salario predicho es 64,000 dólares.
		
		\item Intervalo de confianza del 95\% para $\beta_1$:
		\begin{align}
			\hat{\beta}_1 \pm t_{0.025, 97} \times SE(\hat{\beta}_1) &= 2 \pm 1.985(0.3) \\
			&= 2 \pm 0.5955 \\
			&= (1.4045, 2.5955)
		\end{align}
		
		Con 95\% de confianza, por cada año adicional de experiencia, el salario aumenta entre 1,405 y 2,596 dólares, manteniendo constante el nivel educativo.
	\end{enumerate}
\end{solucion}

\begin{solucion}
	\begin{enumerate}
		\item \textbf{Interpretación de $\hat{\beta}_3 = -3$:} Manteniendo constantes la distancia y el peso, el servicio express reduce el tiempo de entrega en promedio 3 días comparado con el servicio estándar.
		
		\item Prueba F parcial para $X_3$:
		\begin{align}
			F &= \frac{(SSD_{reducido} - SSD_{completo})/(p_{completo} - p_{reducido})}{SSD_{completo}/(n - p_{completo} - 1)} \\
			&= \frac{(78 - 45)/(3 - 2)}{45/(60 - 3 - 1)} = \frac{33/1}{45/56} = \frac{33}{0.8036} \approx 41.07
		\end{align}
		
		Valor crítico: $F_{0.05, 1, 56} \approx 4.01$
		
		Como $F = 41.07 \gg 4.01$, rechazamos $H_0$. La variable $X_3$ es altamente significativa.
		
		\item $R^2$ ajustado:
		
		Modelo completo ($p = 3$):
		\begin{align}
			R^2_{adj} = 1 - \frac{60-1}{60-3-1}(1 - 0.82) = 1 - \frac{59}{56}(0.18) = 1 - 0.1896 = 0.8104
		\end{align}
		
		Modelo reducido ($p = 2$):
		\begin{align}
			R^2_{adj} = 1 - \frac{60-1}{60-2-1}(1 - 0.68) = 1 - \frac{59}{57}(0.32) = 1 - 0.3312 = 0.6688
		\end{align}
		
		\item \textbf{Conclusión:} Sí, definitivamente vale la pena incluir el tipo de servicio porque:
		\begin{itemize}
			\item La prueba F parcial es altamente significativa ($p < 0.001$)
			\item El $R^2$ ajustado mejora sustancialmente (0.810 vs 0.669)
			\item El coeficiente de $X_3$ tiene una interpretación práctica clara y relevante
			\item Reduce significativamente el error del modelo ($SSD$ pasa de 78 a 45)
		\end{itemize}
	\end{enumerate}
\end{solucion}

\begin{solucion}
	\begin{enumerate}
		\item \textbf{Significancia global:} Sí, el modelo es globalmente significativo. El estadístico $F = 35.2$ con $p < 0.001$ indica que al menos uno de los predictores tiene una relación significativa con el rendimiento académico.
		
		\item \textbf{Significancia individual} (con $\alpha = 0.05$):
		\begin{itemize}
			\item $X_1$ (horas de estudio): $p < 0.001 < 0.05$ → \textbf{Significativa}
			\item $X_2$ (asistencia): $p = 0.023 < 0.05$ → \textbf{Significativa}
			\item $X_3$ (recursos en línea): $p = 0.145 > 0.05$ → \textbf{No significativa}
		\end{itemize}
		
		\item \textbf{Acción respecto a $X_3$:} Considerar eliminar $X_3$ del modelo porque:
		\begin{itemize}
			\item No es estadísticamente significativa
			\item No contribuye significativamente al poder predictivo del modelo
			\item Simplificaría el modelo sin perder información importante
		\end{itemize}
		
		Sin embargo, antes de eliminarla, verificar si hay multicolinealidad con otras variables.
		
		\item \textbf{Efecto de eliminar $X_3$:}
		\begin{itemize}
			\item $R^2$: Disminuirá ligeramente (siempre disminuye al eliminar variables)
			\item $R^2_{adj}$: Probablemente aumentará, porque $X_3$ no contribuye significativamente y el $R^2_{adj}$ penaliza la inclusión de variables innecesarias
		\end{itemize}
		
		Esto confirmaría que eliminar $X_3$ mejora el modelo ajustado.
	\end{enumerate}
\end{solucion}

